\chapter*{Resumo}
\pagenumbering{gobble}

% NOTA: o estudo encontra-se em fase de estruturacao metodologica e aprovacao
% etica; ainda nao ha coleta de dados. Por isso, este resumo descreve o
% desenvolvimento do jogo e o protocolo de validacao proposto. Os trechos de
% RESULTADOS e CONCLUSAO estao como PLACEHOLDER e devem ser preenchidos apos a
% coleta e a analise.

O estresse e a ansiedade são queixas frequentes entre jovens adultos e universitários, e intervenções breves, acessíveis e agradáveis para a sua atenuação a curto prazo permanecem um problema em aberto. Este trabalho apresenta o desenvolvimento do \textit{Cozy Pong}, um jogo de realidade virtual (RV) que une mecânicas de tênis de mesa e de jogos de ritmo, ambientado sob a influência de músicas \textit{lo-fi}, e propõe um protocolo para validar empiricamente seu efeito sobre a regulação do estresse e da ansiedade. O jogo foi desenvolvido na \textit{engine} Unity para o dispositivo Meta Quest 3S e organiza a jogabilidade em torno de mapas rítmicos (\textit{beatmaps}) sincronizados à música, com retroalimentação visual, sonora e tátil. Quatro condições operacionalizam a manipulação experimental: o \textit{Cozy Pong} em configuração relaxante, que é a intervenção avaliada; o mesmo jogo em configuração animada, que isola o efeito da caracterização deliberadamente relaxante; a audição da mesma trilha \textit{lo-fi}, sentado e sem jogar, que mantém o estímulo musical constante e isola o efeito do jogo; e uma sessão de meditação guiada no mesmo dispositivo de RV, comparador que representa a prática de relaxamento digital hoje consolidada. O protocolo adota desenho intra-sujeito randomizado e contrabalanceado por quadrado de Williams, com trinta e dois participantes, cada condição precedida de indução padronizada de estresse por tarefa de aritmética mental. O efeito é mensurado pelo cruzamento de medidas psicofisiológicas, a variabilidade da frequência cardíaca (VFC) obtida por cinta peitoral de eletrodos e a atividade eletrodérmica (EDA) obtida por sensor de sinal bruto, com uma bateria psicométrica integralmente composta por escalas de uso livre e com evidências de validade em português brasileiro. A hipótese central é a de que a prática de exercícios leves e ritmados de forma síncrona com músicas relaxantes promove redução mais significativa de estresse e ansiedade do que a audição passiva da mesma música e do que uma configuração mais exigente do próprio jogo, equiparando-se à meditação guiada. \ph{sintetizar os principais resultados: variações nas escalas psicológicas e nos índices fisiológicos por condição.} \ph{enunciar a conclusão principal do estudo após a análise.} Espera-se que o estudo contribua com evidência empírica na intersecção entre computação e saúde a respeito do efeito de intervenções lúdicas em RV sobre a recuperação emocional a curto prazo.

\vspace{1em}
\noindent\textbf{\textit{Palavras-chave}:~} realidade virtual. jogo de ritmo. \textit{exergame}. redução de estresse. ansiedade. música \textit{lo-fi}. variabilidade da frequência cardíaca. atividade eletrodérmica.
